\begin{abstract}
Retail demand forecasting is complicated by heterogeneous data sources, irregular seasonality, intermittent demand for slow-moving items, and the reluctance of business stakeholders to act on predictions they cannot interpret. This paper presents the design and implementation of an end-to-end retail demand forecasting system that trains nineteen forecasting models spanning statistical, machine learning, deep learning, and intermittent-demand families in parallel on a single uploaded dataset, and automatically selects the best-performing model using a robust, multi-criteria composite score rather than a single error metric. The system is implemented as a FastAPI backend and a React single-page frontend, accepts CSV, Excel, and JSON files describing arbitrary time-indexed retail, energy, or transactional data, and performs automated column detection, gap filling, outlier capping, and optional exogenous-feature integration (price and promotional indicators) for the models capable of using them. Beyond numerical accuracy, the system addresses the interpretability gap directly: native feature-importance scores from the tree-based models are surfaced to the user, and a retrieval-augmented large language model layer, grounded in a curated corpus of retail domain knowledge and in the deterministic metrics computed by the pipeline, generates natural-language explanations, a conversational question-answering interface, and a proactive business insights report. Particular engineering attention is given to trustworthiness: forecast uncertainty intervals are derived from each model's own historical cross-validation error rather than an arbitrary fixed percentage, and the language model is explicitly constrained to never contradict the deterministic trend and volatility labels the system itself computes. The system is evaluated on a real retail transaction dataset aggregated to a daily sales series; twelve of the nineteen registered models complete successfully within the automatically inferred data-frequency tier, with the k-nearest-neighbor regressor selected as the best model under the composite score despite extreme gradient boosting achieving a marginally lower root-mean-squared error, illustrating the practical difference between single-metric and multi-criteria model selection. The results and architecture indicate that a broad, automatically-adjudicated model ensemble combined with a grounded natural-language explanation layer is a practical and extensible approach to trustworthy retail forecasting.
\end{abstract}

\begin{IEEEkeywords}
demand forecasting, retail analytics, time series analysis, machine learning, deep learning, ensemble model selection, explainable artificial intelligence, retrieval-augmented generation
\end{IEEEkeywords}
